\lesson{II.2}{Pay for what you want}{Stop choosing gains. Say what an error costs and what thrust costs, and let one equation find the best gains there are.}{1}{lqr}{Part II · Beyond PID}
\label{les:lqr}

\lead{\setupcard{\setuprow{Level}{L1}\setuprow{Controller}{LQR with $q_e = 100$, $q_v = 10$, $r = 10$}\setuprow{Ghost}{the PD of Part~I, $K_p = 10$, $K_d = 7$}\setuprow{Setpoint}{a square wave of $\pm\SI{0.5}{m}$}\setuprow{Goal}{settle within \SI{1.5}{s}, overshoot below 2\,\%, no more peak thrust than the PD}}%
Tuning a PID is haggling with three knobs. This lesson replaces the haggling by a statement of preferences: how much a metre of error costs, how much a metre per second of speed, how much a newton of thrust. A cost function adds them up over the whole flight, and an equation named after Jacopo Riccati\index{Riccati, Jacopo}\index{Riccati equation!in the LQR} returns the gains that make the total as small as it can be. The result is a $K_p$ and a $K_d$ — the same controller as before, chosen by a principle instead of by trial.}

\section{A cost instead of knobs}

The state of the drone, measured from the setpoint, is $\vx = (x, v)$ with $x = y - r$ (\lessonref{les:state}). Any controller produces some motion $\vx(t)$ and some thrust above hover, $u(t) = T - mg$. Price the whole flight:
\begin{equation}\label{eq:lqr-J}
  J = \int_0^\infty \big(q_e\,x^2 + q_v\,v^2 + r\,u^2\big)\,\dd t = \int_0^\infty \big(\vx^\T Q\,\vx + r\,u^2\big)\,\dd t, \qquad Q = \begin{pmatrix} q_e & 0\\ 0 & q_v\end{pmatrix} .
\end{equation}
Each term charges for something one does not want: being away from the setpoint, moving fast, pushing hard. The squares make large deviations disproportionately expensive and do not care about the sign.\tip{Why squares? They are smooth, they punish large errors more than small ones, and they make the answer a linear feedback. Other costs give other — usually nonlinear — controllers.}

\begin{definition}{Linear–quadratic regulator}
For a linear plant $\dot{\vx} = A\vx + B\symbf u$ and the cost $J = \int_0^\infty(\vx^\T Q\vx + \symbf u^\T R\,\symbf u)\,\dd t$ with $Q \succeq 0$ and $R \succ 0$, the \term{linear–quadratic regulator} (LQR) is the controller that makes $J$ smallest from every initial state.
\end{definition}

That such a controller exists, and what it looks like, is the content of one theorem.

\begin{theorem}[The LQR]
\label{thm:lqr-care}
Let $(A, B)$ be controllable and let $Q = C^\T C$ with $(A, C)$ observable. Then the algebraic Riccati equation
\begin{equation}\label{eq:lqr-care}
  A^\T P + PA - PBR^{-1}B^\T P + Q = 0
\end{equation}
has exactly one solution $P \succ 0$. The controller that minimises $J$ is the linear state feedback
\begin{equation}\label{eq:lqr-K}
  \symbf u = -K\vx, \qquad K = R^{-1}B^\T P ,
\end{equation}
the closed loop $\dot{\vx} = (A - BK)\vx$ is asymptotically stable, and the smallest cost is $J_{min} = \vx_0^\T P\vx_0$. (Kálmán, 1960; Anderson and Moore, 1990, ch.~3.)
\end{theorem}

Two remarks before the calculation. The optimal controller is \emph{linear} and \emph{static}: nobody assumed that, it comes out. And the same $K$ is optimal for \emph{every} initial state; the matrix $P$ prices them all at once.\recall{Positive definite: $\vx^\T P\vx > 0$ for every $\vx \ne 0$ (Appendix~G). Here it is the price of starting at $\vx$.}\didyouknow{Riccati studied his equation in 1724, as a problem in pure analysis. It waited 236 years for Kálmán to find what it was for.}

\begin{example}[The Riccati equation of the drone, by hand]
\label{ex:lqr-derive}
Solve \cref{eq:lqr-care} for $A = \begin{psmallmatrix}0 & 1\\ 0 & 0\end{psmallmatrix}$, $B = \begin{psmallmatrix}0\\ 1/m\end{psmallmatrix}$, $Q = \diag(q_e, q_v)$ and $R = r$.

\textbf{Step 1 — the unknowns.} $P = \begin{psmallmatrix}p_1 & p_2\\ p_2 & p_3\end{psmallmatrix}$. With $B^\T P = \tfrac1m(p_2\ \ p_3)$, \cref{eq:lqr-K} gives $K = \tfrac{1}{rm}(p_2\ \ p_3)$.

\textbf{Step 2 — the four terms.} $A^\T P = \begin{psmallmatrix}0 & 0\\ p_1 & p_2\end{psmallmatrix}$, and $PA$ is its transpose. $PBR^{-1}B^\T P = \dfrac{1}{rm^2}\begin{psmallmatrix}p_2^2 & p_2p_3\\ p_2p_3 & p_3^2\end{psmallmatrix}$.

\textbf{Step 3 — three equations.} Entry $(1,1)$: $-p_2^2/(rm^2) + q_e = 0$. Entry $(1,2)$: $p_1 - p_2p_3/(rm^2) = 0$. Entry $(2,2)$: $2p_2 - p_3^2/(rm^2) + q_v = 0$.

\textbf{Step 4 — solve from the top.} $p_2 = m\sqrt{rq_e}$ (the positive root: $P$ must be positive definite). Then $p_3 = m\sqrt{r\,(q_v + 2p_2)}$, and $p_1$ follows from the middle equation.

\textbf{Step 5 — the gains.}
\begin{equation}\label{eq:lqr-gains}
  k_1 = \frac{p_2}{rm} = \sqrt{\frac{q_e}{r}}, \qquad k_2 = \frac{p_3}{rm} = \sqrt{\frac{q_v}{r} + 2m\,k_1} .
\end{equation}
$k_1$ multiplies the position: it is a $K_p$. $k_2$ multiplies the velocity: it is a $K_d$ on the measurement. \emph{The LQR does not replace the PD — it tunes it.}

\textbf{Step 6 — the lesson's numbers} ($m = 1$, $q_e = 100$, $q_v = 10$). For $r = 10$: $k_1 = 3.16$, $k_2 = 2.71$. For $r = 1$: $k_1 = 10$, $k_2 = 5.48$. For $r = 0.1$: $k_1 = 31.6$, $k_2 = 12.8$. The simulator, which solves the discrete version at \SI{250}{Hz}, shows $3.15$ and $2.70$, $9.89$ and $5.44$, $30.8$ and $12.5$.
\end{example}

The formulas read like a sentence. Make thrust cheaper ($r$ smaller) and both gains rise; only ratios of the weights matter, because multiplying $q_e$, $q_v$ and $r$ by the same number leaves \cref{eq:lqr-gains} unchanged. And with $q_v = 0$ the closed loop $ms^2 + k_2s + k_1$ has $\zeta = k_2/(2\sqrt{mk_1}) = \sqrt{2mk_1}/(2\sqrt{mk_1}) = 1/\sqrt2$, whatever $q_e$ and $r$ are: a penalty on position alone always gives $\zeta = 0.71$, the \enquote{best compromise} of \lessonref{les:damper}, which the cost discovers by itself.\tip{To make the LQR calmer, raise $q_v$. To make it faster, lower $r$. To make it stiffer without changing the damping much, raise $q_e$ and $q_v$ together.}

\section{The experiment}

The lesson starts with $r = 10$ and a hand-tuned PD as the ghost. The goal asks for something the ghost cannot do: settle within \SI{1.5}{s} where it needs \SI{2.3}{s}, without overshoot and without asking more of the motors.

\simfig{lqr_runs}{The same step for the PD ghost and three thrust costs. Top: altitude; the grey band is $\pm2$\,\% of the step. Bottom: thrust. With $r = 1$ the LQR has the ghost's $K_p$, hence its peak thrust, and settles \SI{0.85}{s} sooner. With $r = 0.1$ it is faster still, but the thrust runs into the motor limit.}{fig:lqr-runs}

\begin{example}[Why $r = 1$ beats the ghost]
\label{ex:lqr-goal}
Compare the LQR with $r = 1$ with the PD ($K_p = 10$, $K_d = 7$) on the step of \cref{fig:lqr-runs}.

\textbf{Step 1 — the peak thrust.} At the instant of a step only the position term acts: $mg + k_1 \cdot \SI{1}{m}$. Both have $k_1 = K_p = 10$: \SI{19.8}{N} each. The simulator measures 19.7 and \SI{19.8}{N}.

\textbf{Step 2 — the damping.} PD: $\zeta = 7/(2\sqrt{10}) = 1.11$, slightly overdamped, with a slow pole at $-2$. LQR: $\zeta = 5.48/(2\sqrt{10}) = 0.87$, poles $-2.74 \pm 1.58j$.

\textbf{Step 3 — the price.} An overshoot of $e^{-\pi \cdot 0.87/0.5} = 0.4\,\%$ for the LQR: invisible, and well within the goal.

\textbf{Step 4 — the settling.} The overdamped PD crawls in along its slow pole: $\ln 50/2 \approx \SI{2}{s}$ plus the start, \SI{2.33}{s} in the simulator. The LQR's pair decays with $e^{-2.74t}$: \SI{1.48}{s} in the simulator, with 0.09\,\% overshoot. Goal met.

\textbf{Step 5 — what the LQR knew.} Nothing that a good hand-tuner could not have found. But it found it from a statement of preferences, in one step, and — as the next example shows — it found the \emph{best} gains for those preferences.
\end{example}

\section{What \enquote{optimal} buys}

\begin{example}[The cost of a step, for every PD]
\label{ex:lqr-cost}
For the step $\vx_0 = (1, 0)$ and the weights $q_e = 100$, $q_v = 10$, $r = 1$, compute the cost $J$ of the LQR and of three PD tunes.

\textbf{Step 1 — the LQR.} By \cref{thm:lqr-care}, $J_{min} = \vx_0^\T P\vx_0 = p_1$. From Example~\ref{ex:lqr-derive}: $p_2 = 10$, $p_3 = 5.48$, $p_1 = p_2p_3 = 54.8$.

\textbf{Step 2 — any other gain.} For a fixed $K$ the cost is quadratic in $\vx_0$ too: $J = \vx_0^\T P_K\vx_0$, where $P_K$ solves the \emph{Lyapunov} equation $(A - BK)^\T P_K + P_K(A - BK) + Q + K^\T RK = 0$ (along the closed loop, $\frac{\dd}{\dd t}\vx^\T P_K\vx = -(\vx^\T Q\vx + ru^2)$; integrate from 0 to $\infty$).

\textbf{Step 3 — three PDs.} Solving it as in Example~\ref{ex:G-P}: the ghost $(10, 7)$ costs 56.4; the critical $(10, 6.3)$ 55.3; the goal tune of \lessonref{les:state}, $(20, 8)$, 63.8.

\textbf{Step 4 — the picture} (\cref{fig:lqr-cost}, left). Over the whole plane of gains the cost has one minimum, and it is the LQR. The ghost is close — a good hand-tune usually is — but the fast tune of Lesson~II.1 is far: by these weights it pays too much thrust for its speed. \enquote{Optimal} is always optimal \emph{for a cost}; change the weights and the minimum moves.\pitfall[-75pt]{\enquote{Optimal} is not \enquote{best}. It is best for the cost you wrote down — and a cost that forgets the noise, the lag or a limit can make an optimal controller a poor one.}
\end{example}

\simfig{lqr_cost}{Left: the cost of a \SI{1}{m} step for every PD in the plane of $K_p$ and $K_d$, computed from the Lyapunov equation of Example~\ref{ex:lqr-cost}. The LQR sits at the minimum. Right: the LQR gains as the thrust cost $r$ changes, from \cref{eq:lqr-gains} (lines), and the discrete gains the simulator computes at \SI{250}{Hz} (dots).}{fig:lqr-cost}

\subsection{Choosing the weights}
The weights are not physical constants; they are a design language. A common first guess, \emph{Bryson's rule}, takes each weight as one over the square of the largest value one would accept:
\begin{equation}\label{eq:lqr-bryson}
  q_i = \frac{1}{x_{i,max}^2}, \qquad r = \frac{1}{u_{max}^2} .
\end{equation}
For the drone, \enquote{\SI{10}{cm} of error, \SI{1}{m/s} of speed and \SI{5}{N} of extra thrust are acceptable} gives $q_e = 100$, $q_v = 1$, $r = 0.04$, hence $k_1 = 50$, $k_2 = 11.2$ and poles $-5.6 \pm 4.3j$. A starting point, to be adjusted in flight.\pitfall{Large ratios of the weights mean large gains. Every gain the Riccati equation returns must still pass the tests of Part~I: the motor limits, the lag (\lessonref{les:edge}), the noise on the velocity (\lessonref{les:noise}). The cost knows none of them unless they are in the model.}

\screen[0 0 495 998]{lqr}{Lesson II.2 in the app. The info card shows the gains the Riccati equation returned — a $K_p$ and a $K_d$ — and the closed-loop poles; the grey ghost is the PD of Part~I.}{fig:lqr-screen}

\begin{inapp}
\begin{itemize}[leftmargin=1.2em]
  \item The controller is \src{src/control/lqr.ts}. The function \texttt{designLqr} builds the model, discretises it exactly at the controller's period (\texttt{c2d} in \src{src/math/mat.ts}, Appendix~D) and solves the discrete Riccati equation by iteration (\texttt{dlqr} in \src{src/math/riccati.ts}).
  \item The weights are \param{control.lqr.qPos}, \param{qVel} and \param{r}; the switch is \ui{Controller\then Altitude controller\then LQR / LQI}. The gains are recomputed only when a weight, the model or the rate changes.
  \item The info card lists the gains, the closed-loop poles (the eigenvalues of the discrete $A - BK$, converted with $s = \ln z/\Delta t$) and the dominant $\omega_n$ and $\zeta$.
  \item The ghost is flown headlessly when the lesson starts, with the same seed; the goal compares the peak thrust with the ghost's.
\end{itemize}
\end{inapp}

\begin{trythis}
\begin{enumerate}
  \item Lower $r$ from 10 to 1, then to 0.1. Watch the gains, the poles and the thrust. At 0.1 the motors saturate on every step.
  \item With $r = 1$, raise $q_v$ to 100. What happens to $k_2$, to the poles, to the settling time?
  \item Multiply all three weights by ten. Nothing changes — why?
\end{enumerate}
\predict{with $q_v = 0$ and any $q_e$ and $r$, what overshoot will every step have?}
\end{trythis}

\section{The engineer's view: why optimal control scales}

\subsection{Preferences scale; knobs do not}
For two states the LQR is a convenient way to find two gains. Its real value appears with more. The quadrotor of \lessonref{les:cascade} has twelve states and four motors; its full state feedback has 48 gains. Nobody can tune 48 knobs by hand, and a cascade is the classical way around the problem: split the system so that each loop has one or two. The LQR goes the other way. One writes down a $12 \times 12$ diagonal $Q$ and a $4 \times 4$ diagonal $R$ — sixteen numbers with a meaning — and the Riccati equation returns all 48 gains at once, with the couplings between the axes taken into account.\didyouknow{The first great users of the LQR were the aerospace programmes of the 1960s. Today it is the default first design for every new multi-input plant, from satellites to wind turbines.}

\subsection{The guaranteed margins}
Hand-tuned loops are stable if the tuner was careful. LQR loops come with a certificate.

\begin{theorem}[Robustness of the LQR]
\label{thm:lqr-margins}
For a single-input LQR with $R = r > 0$, and the loop broken at the plant input, $L(s) = K(s\Id - A)^{-1}B$:
\begin{enumerate}[label=(\roman*)]
  \item $|1 + L(j\omega)| \ge 1$ at every frequency;
  \item hence the closed loop stays stable if the gain is multiplied by any factor between $\tfrac12$ and $\infty$, and the phase margin is at least $60^\circ$.
\end{enumerate}
(Kálmán, 1964; Safonov and Athans, 1977; Anderson and Moore, 1990, ch.~5.)
\end{theorem}

Statement (i) says that the Nyquist curve never enters the unit disc around $-1$ (Appendix~C); the margins follow by geometry. For the drone with $r = 1$, $L(s) = (k_1 + k_2s)/s^2$; its gain is one at $\omega_c = \SI{5.75}{rad/s}$, where the phase is $-180^\circ + \arctan(5.48 \cdot 5.75/10) = -107.6^\circ$: a margin of $72^\circ$.

The certificate has fine print. It holds where the loop is broken: at the plant's input, for a model that is exact, with the \emph{whole} state measured. The motor lag is outside the model, and it eats into the margin as in \lessonref{les:slow}. And when the state is estimated instead of measured, the guarantee disappears entirely — a famous result of 1978, and the subject of Lesson~III.10.

\begin{example}[Balancing on two wheels]
\label{ex:lqr-segway}
A self-balancing scooter keeps its rider upright by accelerating its wheels. Model the rider as an inverted pendulum with the centre of mass $L = \SI{0.8}{m}$ above the axle: the tilt $\theta$ obeys $\ddot\theta = (g/L)\,\theta - a/L$, where $a$ is the acceleration of the base. Design an LQR with Bryson's rule for $3^\circ$ of tilt, \SI{0.5}{rad/s} of tilt rate and \SI{2}{m/s^2} of acceleration.

\textbf{Step 1 — the plant.} $\vx = (\theta, \omega)$, $A = \begin{psmallmatrix}0 & 1\\ g/L & 0\end{psmallmatrix}$, $B = \begin{psmallmatrix}0\\ -1/L\end{psmallmatrix}$. The eigenvalues of $A$ are $\pm\sqrt{g/L} = \pm3.50$: without control the rider falls with $e^{3.5t}$, doubling the tilt every \SI{0.2}{s}.

\textbf{Step 2 — the weights.} $q_\theta = 1/0.052^2 = 370$ (with $3^\circ = \SI{0.052}{rad}$; Bryson's rule is approximate, take 400), $q_\omega = 1/0.5^2 = 4$, $r = 1/2^2 = 0.25$.

\textbf{Step 3 — solve.} The Riccati equation, solved numerically, gives $K = (-51.0,\ -9.9)$: the base accelerates by \SI{51}{m/s^2} per radian of tilt and \SI{9.9}{m/s^2} per radian per second, \emph{towards} the side the rider leans to.

\textbf{Step 4 — the closed loop.} Poles $-6.2 \pm 3.7j$: stable, $\zeta = 0.86$. The loop crosses unity gain at \SI{12.4}{rad/s} with $67^\circ$ of phase margin, and stays stable if the gain falls to 0.19 of its design value — further than the guaranteed $\tfrac12$.

\textbf{Step 5 — read it.} The plant is unstable, which a PD tuned by feel handles poorly. The LQR does not notice the difference: the same equation, the same certificate. Every self-balancing robot, scooter and rocket on its engine (Lesson~IV.7) begins with a model like this one.\innumbers{An unstable pole at $+3.5$ doubles the tilt every $\ln 2/3.5 = \SI{0.2}{s}$. A human reacts in about the same time — which is why balancing by hand is hard and by computer easy.}
\end{example}

\subsection{In formal terms}

\begin{theorem}[Optimality through completing the square]
\label{thm:lqr-proof}
Under the hypotheses of \cref{thm:lqr-care}, for every input $\symbf u(t)$ that brings the state to zero,
\begin{equation*}
  J = \vx_0^\T P\vx_0 + \int_0^\infty (\symbf u + K\vx)^\T R\,(\symbf u + K\vx)\,\dd t .
\end{equation*}
\end{theorem}
\begin{proof}
Along any trajectory, $\frac{\dd}{\dd t}(\vx^\T P\vx) = \vx^\T(A^\T P + PA)\vx + 2\symbf u^\T B^\T P\vx$. Replace $A^\T P + PA$ by $PBR^{-1}B^\T P - Q$ from \cref{eq:lqr-care} and $B^\T P$ by $RK$: the right-hand side becomes $(\symbf u + K\vx)^\T R(\symbf u + K\vx) - \vx^\T Q\vx - \symbf u^\T R\symbf u$. Integrate from 0 to $\infty$; the left side gives $-\vx_0^\T P\vx_0$, since $\vx \to 0$. Rearrange.
\end{proof}

The proof is the whole theory in one line. The integral is non-negative and is zero exactly when $\symbf u = -K\vx$ at every instant; so no input does better than the LQR, and the LQR achieves $\vx_0^\T P\vx_0$. Existence and uniqueness of the positive definite $P$, and the stability of $A - BK$, are where the hypotheses on controllability and observability are needed (Appendix~I).

The simulator uses the \emph{discrete} version. For $\vx_{k+1} = A_d\vx_k + B_du_k$ and $J = \sum(\vx_k^\T Q_d\vx_k + u_k^\T R_du_k)$ the gain is $K = (R_d + B_d^\T PB_d)^{-1}B_d^\T PA_d$ with $P$ from the discrete Riccati equation $P = Q_d + A_d^\T PA_d - A_d^\T PB_d\,K$. At \SI{250}{Hz} its gains differ from the continuous ones by about one per cent (Example~\ref{ex:lqr-derive}, Step~6); at low rates the discrete design is the right one, because it knows about the hold.

\begin{lookahead}
\begin{itemize}[leftmargin=1.2em]
  \item The LQR has no integral and droops under a wrong model. Lesson~II.3 adds one, as a state.
  \item The LQR needs the whole state. When it is estimated by the Kalman filter of Lesson~II.4 the combination is called LQG, and Lesson~III.10 shows that its margins can be arbitrarily small.
  \item With limits on the thrust the cost cannot simply be minimised by a fixed gain. Model predictive control (Lesson~II.13) minimises the same quadratic cost over a finite horizon, with the limits as constraints, at every sample.
  \item The guaranteed margins are a robustness property; Lessons~III.11–III.14 ask how large a model error any controller can survive. The cost itself returns in Part~IV as the starting point of optimal control — Pontryagin's principle and dynamic programming (IV.1–IV.3), of which the LQR is the one case with a closed-form answer.
\end{itemize}
\end{lookahead}

\begin{remember}
\begin{itemize}
  \item The LQR minimises $\int(\vx^\T Q\vx + \symbf u^\T R\symbf u)\,\dd t$; the answer is the state feedback $\symbf u = -K\vx$, $K = R^{-1}B^\T P$, with $P$ from the Riccati equation.
  \item For the drone: $k_1 = \sqrt{q_e/r}$ (a $K_p$), $k_2 = \sqrt{q_v/r + 2mk_1}$ (a $K_d$). With $q_v = 0$ the damping is always $1/\sqrt2$.
  \item Only ratios of the weights matter. Bryson's rule: one over the square of the acceptable size.
  \item An LQR loop has at least $60^\circ$ of phase margin and a gain margin from $\tfrac12$ to infinity — for the exact model, with the full state measured.
  \item The value of the method is scale: sixteen meaningful weights instead of 48 knobs.
\end{itemize}
\end{remember}

\begin{curiosity}{Balancing on two wheels}
\photo{segway}{A self-balancing scooter of 2006. Gyroscopes and accelerometers measure the tilt; a feedback loop drives the wheels under the rider's centre of mass, hundreds of times per second.}
When the self-balancing scooter was unveiled in 2001, it was sold as a revolution in transport. It was not one — but as a piece of control engineering it is a small masterpiece, and a perfect picture of this lesson.

A person standing on two coaxial wheels is an inverted pendulum, and an inverted pendulum falls: tilted by a degree, it is tilted by two a fifth of a second later (Example~\ref{ex:lqr-segway}). The machine catches the fall by moving the wheels under the centre of mass. Leaning forward makes it accelerate forward, and to keep the lean it must keep accelerating — which is how the rider steers its speed simply by leaning. The controller sees two numbers, the tilt and its rate, and multiplies each by a gain. Choosing those gains well is precisely the problem the LQR solves.

The inverted pendulum on a cart has been the standard laboratory experiment of control courses since the 1960s, because it is the simplest machine that cannot stand without feedback. Two-wheeled robots, rockets balancing on their engines and walking robots are all its descendants.
\end{curiosity}

\begin{exercises}
\exercise[1] With $q_e = 100$, $q_v = 10$ and $m = \SI{1}{kg}$, which thrust cost $r$ gives $k_1 = 20$? What is then $k_2$, and the damping ratio?
\answer{$r = q_e/k_1^2 = 0.25$; $k_2 = \sqrt{40 + 40} = 8.94$; $\zeta = 8.94/(2\sqrt{20}) = 1.0$.}
\exercise[1] Show that multiplying $q_e$, $q_v$ and $r$ by the same factor $c$ leaves the gains unchanged. What happens to $P$ and to $J_{min}$?
\answer{\Cref{eq:lqr-gains} contains only ratios. $P$ and $J_{min}$ are multiplied by $c$.}
\exercise[1] Use Bryson's rule for an error of \SI{20}{cm}, a speed of \SI{2}{m/s} and \SI{10}{N} of extra thrust. Find the gains and the poles.
\answer{$q_e = 25$, $q_v = 0.25$, $r = 0.01$: $k_1 = 50$, $k_2 = \sqrt{25 + 100} = 11.2$; poles $-5.6 \pm 4.3j$ — the same as in the text, because the ratios are the same.}
\exercise[2]\exkind{derive} Complete Example~\ref{ex:lqr-derive}: find $p_1$ and show that $P$ is positive definite. Then compute $J_{min}$ for a step of \SI{0.5}{m}.
\answer{$p_1 = p_2p_3/(rm^2) = m\,k_1k_2 \cdot r$; for the lesson with $r = 1$: $p_1 = 54.8$. Minors: $p_1 > 0$, $p_1p_3 - p_2^2 = 54.8 \cdot 5.48 - 100 = 200 > 0$. $J_{min} = 0.25\,p_1 = 13.7$.}
\exercise[2]\exkind{derive} Show that with $q_v = 0$ the LQR of the drone has $\zeta = 1/\sqrt2$ for every $q_e$ and $r$, and find the overshoot after a step. What does $q_v > 0$ do to $\zeta$?
\answer{$\zeta = k_2/(2\sqrt{mk_1}) = \sqrt{2mk_1}/(2\sqrt{mk_1}) = 1/\sqrt2$; overshoot $e^{-\pi} = 4.3\,\%$. A positive $q_v$ raises $k_2$ and hence $\zeta$.}
\exercise[2]\exkind{derive} Verify the margins of \cref{thm:lqr-margins} for the drone with $r = 1$: find the crossover frequency of $L(s) = (10 + 5.48s)/s^2$, the phase margin, and show that $|1 + L(j\omega)| \ge 1$ at $\omega = \SI{3}{rad/s}$.
\answer{$\omega^4 = 100 + 30\omega^2$: $\omega_c = 5.75$; phase $-180 + 72.4$, margin $72.4^\circ$. At $\omega = 3$: $L = (10 + 16.4j)/(-9) = -1.11 - 1.83j$, $|1 + L| = |-0.11 - 1.83j| = 1.83 \ge 1$.}
\exercise[2]\exkind{fly} In Lesson~II.2 set $r = 1$ and then switch \ui{Model the motor lag} on (\param{control.lqr.lagState}). What do the gains and the poles become, and does the step response change?
\answer{The gains become $(9.89,\ 5.72,\ 0.16)$: $k_2$ rises a little to make up for the lag, and a small third gain feeds back the actual thrust. The poles are $-2.74 \pm 1.59j$ and $-33.2$, the motor. The step hardly changes: at $r = 1$ the loop is far from the edge of \lessonref{les:edge}. The lag state matters at small $r$, where the gains are large.}
\exercise[2]\exkind{code} In \src{src/control/lqr.ts}, \texttt{designLqr} scales the weights by $\Delta t$ before calling \texttt{dlqr}. Why does the discrete cost need that factor to approximate the continuous one? What would change in the gains without it?
\answer{$\sum(\vx_k^\T Q\vx_k + ru_k^2)\Delta t$ approximates $\int(\dots)\,\dd t$. Scaling $Q$ and $R$ by the same factor leaves $K$ unchanged, so the gains would be the same; only $P$, the cost, would be off by $1/\Delta t$.}
\exercise[3]\exkind{derive} The cost of a PD. For the PD $(K_p, K_d)$ with $m = 1$, write $W = Q + rK^\T K$ and solve the Lyapunov equation $(A - BK)^\T P_K + P_K(A - BK) + W = 0$ by hand: show that $p_{12} = w_{11}/(2K_p)$, $p_{22} = (2p_{12} + w_{22})/(2K_d)$ and $p_{11} = K_pp_{22} + K_dp_{12} - w_{12}$. Evaluate the cost of the step $(1, 0)$ for the ghost $(10, 7)$ and compare with the LQR's 54.8.
\answer{The three entries of the matrix equation give the three formulas in turn. For $(10, 7)$ with $q_e = 100$, $q_v = 10$, $r = 1$: $w_{11} = 200$, $w_{12} = 70$, $w_{22} = 59$; $p_{12} = 10$, $p_{22} = 5.64$, $p_{11} = 56.4 + 70 - 70 = 56.4$. The ghost pays 3\,\% more than the optimum.}
\end{exercises}
